% DRAFT (scratch) -- supplement proofs for the clock rule and the absolute bound.
% Not yet in the paper. Framing of the clock rule waits on prechecks/timeout_rule.

\subsection{A Clock in Place of the Counter, and an Absolute Bound}
\label{sec:supp_clock}

\begin{lemma}[Fluid comparison]
\label{lem:fluid}
Let $V(t)$ be the unfinished work at $t$ of a single server of rate $k$ fed the same
work at the same instants, serving whenever it holds work. For every
non-preemptive work-conserving $k$-server policy $Q$ and every $t$,
$V(t) \le U_Q(t) \le V(t) + (k-1)L$.
\end{lemma}

\begin{proof}
Both quantities rise by $C_j$ at the arrival of $j$, so $D = U_Q - V$ changes only
between arrivals, where $U_Q$ falls at the number of busy servers and $V$ falls at
rate $k$ while it is positive. If $D > 0$ then $U_Q > V \ge 0$ and $D$ does not
fall while $V > 0$, and $V$ reaching $0$ leaves $D = U_Q \ge 0$; since $D(0) = 0$,
$D \ge 0$ throughout. If $D > (k-1)L$ then $U_Q > (k-1)L$, so some job waits, since
at most $k-1$ jobs in service hold at most $(k-1)L$; by work conservation all $k$
servers are busy, $U_Q$ falls at rate $k$, and $D$ does not rise. Hence
$D \le (k-1)L$ throughout.
\end{proof}

Write $V_i^-$ for the fluid backlog at $a_i$ of the jobs of rank below $i$, read after
the arrivals at $a_i$: $V_i^- = \max(0, V_{i-1}^- + C_{i-1} - k(a_i - a_{i-1}))$ with
$V_1^- = 0$. Arrivals move $U_Q$ and $V$ by the same amount, so the difference
$U_Q - V$ is the same between any two of the arrivals at one instant as after all of
them. Read it at $a_i$ after the arrivals of rank below $i$ and before $i$'s own: there
$U_Q = R^Q_i$ and $V = V_i^-$, so Lemma~\ref{lem:fluid} gives
$R^Q_i \le V_i^- + (k-1)L$ for every policy $Q$, a clairvoyant one included, since the
lemma is a statement about the schedule actually run.

\begin{proposition}[An absolute bound from the arrivals alone]
\label{prop:absolute}
For every non-preemptive work-conserving policy $P$ and every job $i$,
$k\,W_P[i] \le V_i^- + (k-1)L + \mathrm{In}_i$. In particular
$W_{\mathrm{FCFS}}[i] \le (V_i^- + (k-1)L)/k$, and under Algorithm~1 with cap
$B_{\max}$,
\[
  W_P[i] \;<\; \frac{V_i^- + B_{\max}}{k} + \Bigl(2 - \frac{1}{k}\Bigr)L
         \;\le\; \frac{\sigma + B_{\max}}{k} + \Bigl(2 - \frac{1}{k}\Bigr)L ,
\]
where $\sigma = \max_i (V_i^- + C_i)$ is the smallest depth of a token bucket of
rate $k$ that the arrivals, counted as work, conform to.
\end{proposition}

\begin{proof}
Equation~(busy) of the manuscript gives
$k\,W_P[i] = R^P_i + \mathrm{In}_i - \mathrm{Out}_i - \rho^P_i \le R^P_i + \mathrm{In}_i$,
and $R^P_i \le V_i^- + (k-1)L$. Under first-come first-served $\mathrm{In}_i = 0$.
Under Algorithm~1, $\mathrm{In}_i < B_{\max} + kL$ when $i$ has an overtaker, by
Theorem~3 of the manuscript, and $\mathrm{In}_i = 0$ otherwise. For the last step,
$V_i^- \le \sigma - C_i < \sigma$. That $\sigma$ is the least conforming depth is the
standard reading of the Lindley recursion: $V_i^- + C_i$ is the largest amount by
which the work arriving in a window ending at $a_i$ exceeds $k$ times the window's
length.
\end{proof}

At $k = 1$ the fluid server is the queue itself, $V_i^- = W_{\mathrm{FCFS}}[i]$, and the
proposition is Theorem~3 of the manuscript again. For $k \ge 2$ it is not a
consequence of Theorem~3: going through $W_{\mathrm{FCFS}}$ would add the
$2(k-1)L$ of Theorem~1 a second time.

\begin{corollary}[A cap from an absolute promise]
\label{cor:absolute}
An operator who promises that no job waits longer than $D$, and who expects the
arrivals of the coming period to conform to a rate-$k$ bucket of depth $\sigma$, sets
$B_{\max} = kD - \sigma - (2k-1)L$ when that is non-negative.
\end{corollary}

\begin{proposition}[A clock keeps the same promise]
\label{prop:clock}
Let Timeout($\theta$) serve, at each dispatch epoch, the oldest waiting job if it has
waited at least $\theta$, and the base policy's choice otherwise. For every base
policy, every input and every job $i$ with an overtaker,
$\mathrm{In}_i < k(\theta + L)$, hence
\[
  W_P[i] \;<\; W_{\mathrm{FCFS}}[i] + \theta + (3 - 2/k)L,
  \qquad
  W_P[i] \;<\; V_i^-/k + \theta + (2 - 1/k)L .
\]
\end{proposition}

\begin{proof}
From $a_i + \theta$ on, $i$ has waited $\theta$, so the oldest waiting job has too and
every dispatch until $s_i$ serves a job older than $i$. Every overtaker of $i$ therefore
starts on some server in $[a_i, a_i + \theta)$ and ends before its start plus $L$. The
overtakers that share a server run one after another inside
$[a_i,\, a_i + \theta + L)$, so they hold that server for less than $\theta + L$, and
summing over the $k$ servers gives $\mathrm{In}_i < k(\theta + L)$. Corollary~1 of the
manuscript turns this into the first bound and Proposition~\ref{prop:absolute}'s first
inequality into the second. A job with no overtaker has excess at most $(2-2/k)L$.
\end{proof}

The rule is Algorithm~1 with a budget that is unbounded until age $\theta$ and zero from
then on. Its fired set is a prefix of the queue in arrival order, so the head is
always its minimum-rank member. Read as a charge, the clock bills a waiting job for
every unit of work the $k$ servers execute while it waits, $k(t - a_q)$, and the guard
bills it only for completed work that arrived after it. Firing on either rule keeps
both promises at once, since each proof reads only its own condition; with the clock's
half computed from arrival times alone, the promise
$W < W_{\mathrm{FCFS}} + \theta + (3-2/k)L$ survives completion reports that are late
or lost.
